\hypertarget{appendix-275ec8847057}{}\label{referencestart:appendix-275ec8847057}
\section{Reversed-Z depth contracts}\label{appendix-275ec8847057-reversed-z-depth-contracts}

A depth buffer stores an ordering coordinate, not Euclidean distance. Reversed-Z changes the mapping so near geometry receives larger depth values and far geometry approaches zero. The precision advantage comes from combining the reciprocal structure of perspective depth with floating-point numbers whose representable values are denser near zero.

The entire renderer has to agree on the convention. Reversing only the projection matrix while keeping old clear/comparison assumptions creates a system that is internally contradictory.

\subsection{Contract checklist}\label{appendix-275ec8847057-contract-checklist}

For a conventional zero-to-one API depth range, a common reversed-Z setup uses:

\begin{itemize}
\tightlist
\item
  near plane mapped near 1;
\item
  far distance mapped toward 0;
\item
  depth buffer cleared to 0;
\item
  comparison set to GREATER or GREATER\_EQUAL;
\item
  hierarchical depth/occlusion reductions defined consistently with ``larger is nearer'';
\item
  screen-space reconstruction derived from the exact projection convention.
\end{itemize}

Infinite-far reversed projection can remove an explicit far-plane distance from the perspective matrix, but frustum/culling policy still needs a practical world-space bound when the application cannot process infinite geometry.

\subsection{Why floating point helps}\label{appendix-275ec8847057-why-floating-point-helps}

Perspective depth is approximately affine in \(1/z\). Conventional mapping spends many floating-point representable values near depth one, where far geometry is crowded. Reversed-Z places far geometry near zero, where floating-point exponent scaling gives much more useful relative spacing.

With integer/fixed-point depth, reversal alone does not create the same dramatic distribution advantage; the floating-point representation is part of the argument.

\subsection{Reconstruction must use the actual matrix}\label{appendix-275ec8847057-reconstruction-must-use-the-actual-matrix}

Do not paste a ``linearize depth'' formula from another API. Different clip-space ranges, handedness, finite/infinite far planes and reversed mappings change the coefficients.

A robust reconstruction derives camera-space position from the same projection matrix or stores explicit projection parameters and tests known near/far points.

\subsection{Occlusion and temporal systems}\label{appendix-275ec8847057-occlusion-and-temporal-systems}

Hierarchical Z, SSAO, SSR, clustered depth slicing, picking and temporal rejection all consume depth meaning. Every reduction or comparison must be audited when the convention changes. A min-reduction hierarchy written for standard depth may need a max-reduction interpretation under reversed-Z.

\subsection{Tests}\label{appendix-275ec8847057-tests}

Project known camera-space depths, confirm monotonic ordering, quantize them to the real depth format, reconstruct them, and measure relative world-space error over logarithmically spaced distances. Also test clear-value behavior with an empty frame and a near geometry sample.

Teach and reconstruct: \hyperref[03-visual-computing-i-ch012-give-depth-an-explicit-meaning]{Continuous Notes/University - Worlds, Images and Inference - Graphics, Vision and Photography\#Give depth an explicit meaning}.

Related: \href{../07-Vision-XR-and-Playable-Worlds/main.pdf\#nameddest=appendix-1bc40b08db4a}{Hardware depth-buffer linearization (Vision, XR, and Playable Worlds, Reference notes: Hardware depth-buffer linearization)}, \href{../03-Geometry-Cameras-and-Image-Formation/main.pdf\#nameddest=appendix-2e7e9fd28a2c}{Rendering history validity (Geometry, Cameras, and Image Formation, Reference notes: Rendering history validity)}.
